\documentclass[11pt]{article}

\usepackage[margin=1in]{geometry}
\usepackage{amsmath,amssymb,amsthm,mathtools}
\usepackage{booktabs,array}
\usepackage{microtype}
\usepackage{enumitem}
\usepackage{xcolor}
\usepackage[hidelinks]{hyperref}

% DRAFT (2026-07-12). Manuscript surfacing of the null-edge cosmological-
% constant program (rungs L1-L6 of Null_Edge_Cosmological_Constant_2026-07-09.md).
% Substantive where kernel theorems back it; outline/conjecture where marked.
% RELEASE GATES before this leaves draft: (i) DONE 2026-07-12 - the Lambda
% capstones are guard-pinned in PhysicsSM/Draft/NullEdge/LambdaCosmologyAxiomGuard
% (~44 pins, 11 modules, build green); (ii) named authors; (iii) verify every
% [import] against primary sources.

\newcommand{\Lam}{\Lambda}
\newcommand{\ii}{\mathrm{i}}
\definecolor{kcol}{RGB}{0,92,74}
\definecolor{icol}{RGB}{0,64,128}
\definecolor{ccol}{RGB}{145,76,0}
\definecolor{scol}{RGB}{120,120,120}
% Claim calculus (project convention):
\newcommand{\Kernel}{\textcolor{kcol}{\textsc{M}}}       % kernel-checked in repo
\newcommand{\Import}{\textcolor{icol}{\textsc{[import]}}} % external result relied on
\newcommand{\Conj}{\textcolor{ccol}{\textsc{[C]}}}        % pre-registered conjecture
\newcommand{\Spec}{\textcolor{scol}{\textsc{[spec]}}}     % honest speculation
\newcommand{\lm}[1]{\texttt{#1}}

\newtheorem{theorem}{Theorem}
\newtheorem{proposition}{Proposition}
\newtheorem{remark}{Remark}

\title{The Cosmological Constant as the Count-Conjugate of a Finite
Null-Edge Spacetime:\\ a kernel-checked everpresent-$\Lam$ core and its
pre-registered dichotomy}
\author{[named authors --- release gate]}
\date{July 12, 2026 --- working draft (partial outline)}

\begin{document}
\maketitle

\begin{abstract}
We report a Lean~4 formalization, over Mathlib, of the arithmetic core of a
finite ``null-edge'' account of the cosmological constant $\Lam$.  The value
proposition is exactness and machine-verification of the \emph{structural}
claims, not a first-principles derivation of the observed value.  Three
independently landed, kernel-checked (\Kernel) results organize the account.
(1) \emph{Magnitude dissolution, structurally}: in the finite spectral-action
grading the coefficient of $\operatorname{tr}(\mathbf 1)$ is invariant under
\emph{every} deformation of the Dirac operator, while the order-2 and order-4
traces provably do change; only count statistics can touch the order-0
coefficient.  (This grading invariance is the finite avatar of a classical
heat-kernel fact; the physical statement that no matter or gauge channel
renormalizes $\Lam$ rests on the finite unimodular trade of Section~4, not on
the grading alone --- see the Section~4 remark.)
(2) \emph{The everpresent scaling law}: with the discreteness (Poisson) input
$\langle\delta N^2\rangle = V$ the normalized $\Lam = \delta N/V$ has RMS
$1/\sqrt V$, routed through the framework's own extensive edge count rather than
an abstract volume.  (3) \emph{A pre-registered dichotomy with a sharp
exponent}: the count variance $\operatorname{Var}(N)$ is a finite susceptibility
whose scaling exponent $\alpha$ decides everpresent survival --- the Poisson
branch ($\alpha=1$) gives the everpresent exponent $-\tfrac12$ and predicts
fluctuating dark energy at $1/\sqrt V$, while any hyperuniform branch
($\alpha<1$, the constraint-correlated case realized by Coulomb systems)
suppresses $\Lam$ below $1/\sqrt V$ and kills the everpresent identification.
This dichotomy is now a \Kernel{} theorem on the framework's own fermionic states:
via Wick's number-variance identity, an explicit projection (Fermi-sea) kernel
realizes a genuinely sub-extensive, \emph{unbounded} variance ($\alpha=\tfrac12$),
so the fork is proved and only the physical count-identification remains open.
We additionally give an \emph{illustrative} finite Fourier/uncertainty picture of
the $\Lam$--volume conjugacy on a small register, and pre-register the canonical
finite Henneaux--Teitelboim statement ($\{\Lam,V\}=1$, $\dot\Lam=0$ on shell) as
its proper successor.  Every trust boundary is stated at its
point of use.  What is \emph{not} claimed: a derivation of the value
$10^{-122}$ (the 4-volume $V$ is an input), of the sign, or of the stochastic
dynamics.  As of this draft the capstones are kernel-clean \emph{and}
build-guarded: a dedicated axiom guard (\lm{LambdaCosmologyAxiomGuard}) pins the
footprint of every cited theorem on every build (Section~7).
\end{abstract}

\paragraph{Claim calculus.}  \Kernel{} = a theorem kernel-checked in this
repository (axiom footprint $\{\texttt{propext},\texttt{Classical.choice},
\texttt{Quot.sound}\}$; every module cited below is \texttt{native\_decide}-free).
\Import{} = an external result relied on but not proved here.  \Conj{} = a
pre-registered conjecture with a stated kill condition.  \Spec{} = honest
speculation.  Unmarked prose is interpretation.

\section{The problem, in two halves}

\emph{Magnitude.}  Continuum QFT's naive vacuum zero-point density is
$\sim M_{\mathrm{Pl}}^4$, about $120$ orders above the observed
$\rho_\Lam\sim(10^{-3}\,\mathrm{eV})^4\sim10^{-122}$ in Planck units.
\emph{Coincidence.}  Why is $\rho_\Lam$ comparable to the matter density
\emph{now}, when the two scale differently with expansion?  A finite theory has
no divergent zero-point sum, so the magnitude problem cannot appear in its
continuum form; but finiteness alone does not \emph{predict} smallness, and
matter loops still renormalize the vacuum energy (Weinberg 1989 \Import).  The
framework's genuine leverage is (i) a \emph{structural} reason no channel can
feed the order-0 coefficient (Section~4), and (ii) a specific, falsifiable
handle on the coincidence half (Sections~3,5).

\section{The finite spectral grading and the order-0 coefficient}

In the spectral-action route (Chamseddine--Connes--Marcolli \Import) the action
is one functional whose expansion gives a cosmological/volume term (order~0), a
gravity term (order~2), and matter terms (order~4).  Thus $\Lam$ is the
order-0 coefficient $a_0\operatorname{tr}(\mathbf 1)$ --- a finite, computable
number, and a graded piece of the same functional yielding gravity and matter.

The finite avatar makes this exact.  In \lm{LambdaMomentHierarchy} the graded
pieces are explicit and the following hold \Kernel:
\begin{itemize}[leftmargin=1.6em,itemsep=1pt]
\item \lm{moment\_hierarchy}: the order-$0/2/4$ decomposition of the finite
  functional into $(a_0\operatorname{tr}\mathbf1,\ \text{gravity},\ \text{matter})$;
\item \lm{order0\_deformation\_invariant}: $\operatorname{order0}(a_0,D+P)=
  \operatorname{order0}(a_0,D)$ for \emph{any} perturbation $P$;
\item \lm{only\_count\_touches\_lambda} and \lm{hierarchy\_verdict}: the order-0
  coefficient is blind to every channel deformation, the higher orders are not.
\end{itemize}

\section{The main handle: everpresent $\Lam$}

\paragraph{Mechanism (Sorkin; Ahmed--Dodelson--Greene--Sorkin \Import).}
In unimodular gravity $\Lam$ is conjugate to the spacetime 4-volume $V$; in a
discrete spacetime the element count $N\sim V$, a Poisson fluctuation gives
$\delta N\sim\sqrt V$, and conjugacy makes $\Lam$ fluctuate about zero with
$|\Lam|\sim\delta N/V\sim1/\sqrt V$.  For the present horizon $V\sim10^{244}$
Planck units, so $1/\sqrt V\sim10^{-122}$ --- the observed magnitude, untuned
\Import --- and since $\Lam\sim1/\sqrt V$ tracks the ambient density at every
epoch, the coincidence is dissolved.

\paragraph{Kernel core.}  The scaling arithmetic is machine-checked, first
abstractly and then through the framework's own primitive:
\begin{itemize}[leftmargin=1.6em,itemsep=1pt]
\item \lm{everpresentLambda\_secondMoment\_eq\_inv\_volume},
  \lm{everpresentLambda\_rms\_eq\_inv\_sqrt\_volume} \Kernel{}: given
  $\langle\delta N^2\rangle = V$, $\Lam=\delta N/V$ has second moment $1/V$ and
  RMS $1/\sqrt V$ (\lm{NullEdgeP9EverpresentLambdaScaling}).
\item \lm{edgecount\_extensive}, \lm{edgecount\_mono},
  \lm{lambda\_rms\_eq\_inv\_sqrt\_count}, \lm{everpresent\_verdict} \Kernel{}:
  the pierced-null-edge count $N$ is additive over disjoint regions and
  monotone, and the everpresent RMS holds with $N$ \emph{the edge count}
  (\lm{LambdaEdgeCount}) --- the handle is native, not borrowed via an abstract
  $V$.
\item \lm{uniformSecondMoment\_antitone},
  \lm{uniformSuppression\_below\_everpresent} \Kernel{}: a uniform suppression
  with residual $A^2/N$ beats the $\sqrt N$ fluctuation iff the scale is
  sub-extensive (\lm{NullEdgeP9EverpresentLambdaTension}) --- the comparison
  machinery for Section~5.
\end{itemize}

\paragraph{Honest scope.}  \Kernel{} is the \emph{arithmetic of the scaling}.
The inputs remain \Import/\Conj: the $\Lam$--$V$ conjugacy (unimodular gravity),
the Poisson second moment $\langle\delta N^2\rangle=V$ (Section~5 situates this),
and the identification of the edge count with the 4-volume.

\section{Structural magnitude dissolution (the headline theorem)}

The order-0 blindness of Section~2 upgrades the ``finiteness removes the
divergence'' story to a structural statement: matter physics \emph{cannot}
renormalize $\Lam$.  In \lm{LambdaUnimodular} and the capstones:
\begin{itemize}[leftmargin=1.6em,itemsep=1pt]
\item \lm{trace\_one\_eq\_dim} ($\operatorname{tr}\mathbf1 = n$),
  \lm{trace\_channel\_blind}, \lm{order0\_operator\_blind} \Kernel: the order-0
  coefficient is invariant under channel and operator deformations.
\item \lm{multiplier\_field\_equation}, \lm{vacuum\_shift\_is\_gauge},
  \lm{gauge\_solution\_map} \Kernel: the finite unimodular trade --- under a
  count constraint a shift of the order-0 coefficient is a gauge move on the
  constraint surface (unobservable), the multiplier enters the field equation as
  $+\Lam\,\mathbf1$ (Jacobson's integration constant, finitely), and the
  observable residue is the count's deviation from its constrained mean.
\item \lm{lambda\_only\_count\_can\_move\_order0} \Kernel{} (\lm{LambdaMagnitudeCapstone}):
  bundles it --- $\operatorname{order0}(D+P)=\operatorname{order0}(D)$ and
  $\Lam$ invariant under vacuum shifts, \emph{while}
  $\operatorname{tr}((D+P)^2)\neq\operatorname{tr}(D^2)$ and the order-4 trace
  likewise moves.  Matter physics moves the gravity and matter sectors and
  leaves $\Lam$ untouched; only count statistics can move it.
\item \lm{lambda\_magnitude\_capstone}, \lm{lambda\_nonvacuity\_witnesses} \Kernel:
  the magnitude verdict with explicit non-vacuity witnesses.
\end{itemize}

\begin{remark}[what this is and is not, and where the load-bearing wall really is]
The order-0 invariance is the finite avatar of a \emph{classical} heat-kernel
fact --- the leading Seeley--DeWitt coefficient $a_0$ depends only on
$\text{rank}\times\text{volume}$, not on the operator \Import{} --- so
$\operatorname{order0}$ not mentioning $D$ makes \lm{order0\_deformation\_invariant}
close to definitional (it is the structural encoding of the grading, best read as
the Classical-analogue it is).  Consequently the grading theorem \emph{alone} does
not carry the physical claim ``matter cannot renormalize $\Lam$'': as Section~1's
Weinberg concession already admits, a Higgs-sector vacuum shift sits in the
order-4 trace and still gravitates as an effective $\Lam$.  What neutralizes it is
the finite \emph{unimodular trade} --- \lm{vacuum\_shift\_is\_gauge} together with
the imported $\Lam$--$V$ conjugacy (Section~3): under a fixed-count constraint a
uniform $c\,\mathbf1$ shift is a gauge move on the constraint surface, hence
unobservable.  That trade, not the grading, is the load-bearing wall, and the
paper should be read that way.  This dissolves the magnitude problem
\emph{structurally} for the finite functional but does not derive the value
($\operatorname{tr}\mathbf1=(\text{edge count})\times(\text{fibre dimension})$,
$V$ an input).  The honest finite sequestering core still owes two keystones, both
landable from \lm{LambdaUnimodular} (\emph{pre-registered, not yet theorems}):
(i) an \emph{ensemble} shift-invariance --- at fixed count the partition function
and every observable are invariant under a uniform $c\,\mathbf1$ shift, upgrading
the trade from classical field equation to statistical mechanics; and (ii) a
\emph{local-shift decomposition} --- an inhomogeneous $c(x)\mathbf1$ splits into a
gauged uniform mode and a residue provably routed into the order-2/4 observables.
\end{remark}

\section{The pre-registered dichotomy and the exponent prediction}

The Poisson input need not be assumed: it is a statement about the count
variance, which the finite ensemble computes.

\paragraph{$\Lam$ as a susceptibility.}  \lm{LambdaSusceptibility} \Kernel:
$\lm{expect\_Ncount}$, $\lm{var\_count}$ give $\operatorname{Var}(N)$ for
independent occupancies, and $\lm{bernoulli\_bound}$ proves
$\operatorname{Var}(N)=\sum_i p_i(1-p_i)\le\langle N\rangle$ (equality in the
sparse/Poisson limit).  Everpresent becomes an \emph{upper-bound theorem for an
ideal edge gas}: $\Lam_{\mathrm{rms}}\le1/\sqrt{\langle N\rangle}$, with
$\Lam_{\mathrm{rms}}=\sqrt{\chi_N}/\langle N\rangle$ the count susceptibility.

\paragraph{The fork \Conj (the distinctive question).}  Do the framework's own
constraints correlate the count?  Constraint-induced correlations can make
number fluctuations sub-extensive --- the physical precedent is that Coulomb
systems are hyperuniform (charge fluctuations scale with surface, not volume)
\Import.  \lm{LambdaCountDichotomy} \Kernel{} formalizes both branches on finite
witnesses: \lm{free\_variance\_extensive}/\lm{free\_extensive} (independent
edges, $\operatorname{Var}\sim N$, Poisson) versus \lm{constrained\_variance\_hard}/
\lm{hard\_subextensive} and \lm{soft\_variance}/\lm{soft\_subextensive} (a
constrained witness with strictly sub-extensive variance).  What the kernel
currently proves is that both branches are \emph{possible} on witnesses --- which
nobody doubted; the answer, not the possibility, is the paper, and it is
computable because the counted objects are fermionic.

\paragraph{Resolving the fork on the framework's own fermionic states \Conj{}
(the decisive upgrade).}  For any free-fermion/Gaussian (CAR) state with
correlation matrix $K$ ($0\le K\le1$), Wick's theorem gives the regional number
variance in one line,
$\operatorname{Var}(N_A)=\operatorname{tr}K_A-\operatorname{tr}K_A^2$; the diagonal
case $K=\operatorname{diag}p$ recovers \lm{var\_count} above (the extensive
branch).  A \emph{projection} kernel $K^2=K$ --- a Fermi-sea, i.e.\
\emph{kinematic}, count --- instead makes $\operatorname{Var}(N_A)$ count only
boundary-crossing modes.  An explicit finite family ($b$ fully-occupied bulk sites
plus $k$ boundary bonds) has $\operatorname{Var}=k/4$ with region size tunable to
$k^2$: a genuinely \emph{sub-extensive and unbounded} variance, $\alpha=\tfrac12$
--- the non-degenerate hyperuniform witness the kill branch needs, on a bona fide
kernel rather than a degenerate toy.  The fork then sharpens into a
\emph{dichotomy}: \emph{everpresent $\Lam$ survives iff the $\Lam$-conjugate count
is the thermal/diagonal count rather than the kinematic/projection
(vacuum-weighted) count.}  Two identifiable results tension it from both sides:
the framework's own vacuum (finite Sorkin--Johnston / discrete free-fermion) reading
is generically sub-extensive (Gioev--Klich regime) and would fire the kill, while a
Lorentz-invariance leg (Bombelli--Henson--Sorkin: Poisson is the essentially unique
Lorentz-invariant discretization \Import) pins the \emph{kinematic sprinkling} count
at $\alpha=1$.  The \emph{mathematical} dichotomy is now \Kernel{} (\lm{LambdaFermionicFork},
guard-pinned): the fermionic number-variance identity, an explicit boundary-bond
\emph{projection} kernel ($K^2=K$, kernel-checked --- a genuine Fermi-sea
correlation matrix, \lm{bondProj\_isProjection}) with $\operatorname{Var}=k/4$
exactly (\lm{bondProj\_numberVariance}) and region size tunable to $k^2$
(\lm{fork\_subextensive}: $\alpha=\tfrac12$, sub-extensive \emph{and} unbounded),
and the two-branch verdict (\lm{fermionic\_fork\_verdict}).  So the fork is a
theorem, not merely a possibility: a diagonal/thermal count is extensive and
everpresent survives, while a projection/kinematic count can be genuinely
sub-extensive and everpresent is killed --- on a bona fide kernel, not a
degenerate toy.  What remains \Conj{} is only the \emph{physical}
count-identification: which count the framework's own dynamics makes
$\Lam$-conjugate (bare vs vacuum-weighted).  Either resolution publishes; a
hyperuniform no-go for everpresent $\Lam$ in constrained fermionic discreta would
be the higher-impact outcome.

\paragraph{The exponent, exactly.}  \lm{LambdaExponentFork} \Kernel{} turns the
fork into a single exponent.  With $\operatorname{Var}(N)\sim N^\alpha$ the
$\Lam$ exponent is
\[
  \lm{lamExp}(\alpha)=\tfrac{\alpha}{2}-1
  \qquad(\lm{lamExp\_closed}),
\]
so \lm{everpresent\_value}: $\lm{lamExp}(1)=-\tfrac12$ (Poisson);
\lm{hyperuniform\_faster}: $\alpha<1\Rightarrow$ exponent $<-\tfrac12$
(suppressed below $1/\sqrt V$; e.g.\ area scaling $\alpha=\tfrac34$ gives
$\sim10^{-152}$); \lm{superextensive\_slower}; and \lm{fork\_iff}:
$\lm{lamExp}(\alpha)=-\tfrac12\iff\alpha=1$.  In program discipline the
prediction is an \emph{exponent}, not a scale: a measured deviation of the
$\Lam$-variance exponent from $-\tfrac12$ measures hyperuniformity.

\paragraph{Kill condition \Conj.}  If the framework's Gauss/gauge constraints
force $\alpha<1$ on physical witnesses, the everpresent identification of
observed dark energy dies in this framework and $\Lam$ demotes to an unexplained
integration constant.  Which count is conjugate to $\Lam$ (bare edge count vs
decorated/soldered count) is decidable on finite witnesses; either answer is
publishable.

\section{A finite Fourier/uncertainty picture of the conjugacy (illustrative)}

\lm{LambdaConjugacy} \Kernel{} builds a volume register and a $\Lam$ partner by a
finite DFT over $\mathbb Z/4$ and proves a support-uncertainty relation:
\lm{delta\_maps\_to\_uniform} and \lm{uniform\_maps\_to\_delta} (a sharp count is a
flat $\Lam$ and conversely), \lm{support\_uncertainty} (a Donoho--Stark support
bound), \lm{delta\_saturates} (the extremal register saturates it).  We are careful
not to oversell this: the $\mathbb Z/4$ support-uncertainty relation is an
\emph{illustrative} finite avatar of ``$\Lam$ conjugate to volume'', a
metaphor-made-exact on one small register, \emph{not} the canonical conjugacy
statement.  The dictionary it suggests --- mass $=$ residual mixedness of hidden
null \emph{directions}, $\Lam$ $=$ residual phase noise of hidden null
\emph{volumes} (the same trace-out-the-hidden-register operation on two registers)
--- is interpretation, not theorem.  Two successors make ``native'' literal:
(i) the general-$N$ Donoho--Stark support bound
$N\le|\operatorname{supp}f|\cdot|\operatorname{supp}\widehat f|$ over $\mathbb Z/N$,
now \Kernel{} for \emph{all} $N$ (\lm{LambdaUncertaintyGeneralN.support\_uncertainty},
via a machine-checked Plancherel and Cauchy--Schwarz on Mathlib's \lm{ZMod.dft};
guard-pinned), which retires the ``$\mathbb Z/4$ witness only'' scope of this
section --- its extremizer classification is the remaining piece; and (ii), the
real prize, a finite Henneaux--Teitelboim statement --- a constrained Hamiltonian minisuperspace in
which $\Lam$ is the momentum conjugate to total $4$-volume, $\{\Lam,V\}=1$,
$\dot\Lam=0$ on shell, with \lm{vacuum\_shift\_is\_gauge} as a corollary.  To our
knowledge unimodular gravity's $\Lam$-as-integration-constant has not been
machine-checked even in a minisuperspace avatar; that theorem, not the DFT witness,
is what would make the conjugacy native.  (Both are \Conj{}/pre-registered.)

\section{Verification status, and the one honest gap that remains}

Every module cited is kernel-clean (\texttt{native\_decide}-free; footprint
$\{\texttt{propext},\texttt{Classical.choice},\texttt{Quot.sound}\}$) and
imported in the draft build root.
\begin{enumerate}[leftmargin=1.6em,itemsep=1pt]
\item \textbf{Build-enforced axiom guard: in place.}  Every headline theorem
  cited above is now pinned by a \lm{\#guard\_msgs}/\lm{\#print axioms} block in
  \lm{PhysicsSM/Draft/NullEdge/LambdaCosmologyAxiomGuard.lean} (about forty pins
  across the eleven $\Lam$ modules; build green).  The kernel-clean footprint is
  therefore re-checked on \emph{every} build: should any dependency ever regress
  --- a \texttt{sorry}, an \texttt{axiom}, or a \texttt{native\_decide} entering
  a transitive proof --- the recorded axiom string stops matching and the guard
  module fails to build.  A single pin is instructive: \lm{edgecount\_mono} is
  guarded at the sharper footprint $\{\texttt{propext},\texttt{Quot.sound}\}$
  (it needs no choice), so even a spurious \texttt{Classical.choice} would trip
  the guard.  (Pattern: the sibling \lm{OvernightTheoryAxiomGuard}.)
\item \textbf{This document is the surfacing (the remaining gate).}  Until now
  the results lived only in draft modules; this manuscript is rung L6.  It
  requires named authors and a primary-source pass on every \Import{} before
  leaving draft.
\end{enumerate}

\section{The event horizon (what is \emph{not} claimed)}

No derivation of $V$, hence not of the value $10^{-122}$: $V$ (the age/size of
the universe) is an input.  No derivation of the \emph{sign} (the
positive-sector/Krein bias is a one-line \Spec{} question at most).  No claim on
the stochastic \emph{dynamics} (sequential growth is causal-set theory's, not
ours).  The claim to defend is exactly: the framework \emph{dissolves the
magnitude divergence structurally} (Sections~2,4), \emph{explains the
coincidence conditionally} on the Poisson branch (Section~5), and derives
neither the value nor the sign.

\section{Predictions (conditional on the Poisson branch) \Conj}

\begin{enumerate}[leftmargin=1.6em,itemsep=1pt]
\item $\langle\Lam\rangle\approx0$, $\Lam_{\mathrm{rms}}\sim V^{-1/2}$; the
  observed positive value is a realization, the sign undetermined.
\item Coincidence is \emph{permanent}: $|\rho_\Lam|\sim\rho_{\mathrm{crit}}$ at
  every epoch, tiny because the universe is old and large.
\item Deviations are \emph{horizon-scale} (low-frequency drift), not a
  clustering local fluid.
\item CPL $(w_0,w_a)$ fits may \emph{mislead}: stochastic $\Lam(t)$ can fake
  smooth evolution; the object is volume noise, not a rolling scalar.
\item No exact future de~Sitter equilibrium: $\Lam_{\mathrm{rms}}\sim t^{-2}$
  keeps tracking; long-term acceleration may be intermittent.
\end{enumerate}

\noindent\emph{Inherited caveat (out of claimed scope).}  Predictions~1 and~5 lean
on temporal structure this paper does not supply: a zero-mean $\Lam$ that is
observed \emph{persistently} positive over $\sim$Hubble times requires long
correlation times in the stochastic history, which is Zwane--Afshordi--Sorkin
territory (the dynamics of everpresent $\Lam$), not the finite kinematic core
formalized here.  We flag rather than assume it; a genuine prediction of the sign's
persistence would need that dynamical input.

\paragraph{Observational posture \Import{} (verify before citing).}
Planck 2018 sets rigid $\Lam$CDM as baseline; DESI DR2 BAO$+$CMB report a
time-evolving equation of state at $\sim3.1\sigma$ ($2.8$--$4.2\sigma$ with
supernovae) --- \emph{hints}, not a replacement.  The framework inherits a
falsifiable phenomenology (fluctuating, sign-changing, $1/\sqrt V$-tracking dark
energy) \emph{conditional on the Poisson branch}.  Live tension:
Zwane--Afshordi--Sorkin (everpresent models can fit as well as $\Lam$CDM and
relieve low-$z$ tensions) vs the Aspects-II finding (CMB-favored fluctuation
magnitudes smaller than typical everpresent models).  A confirmed rigid
$w=-1$ with tight fluctuation bounds is the kill; a confirmed deviation is
friendly but not a confirmation.

\section*{Landed-theorem table}

\begin{center}\small
\begin{tabular}{@{}p{0.46\linewidth}p{0.34\linewidth}c@{}}
\toprule
Statement & Lean declaration & Grade\\
\midrule
order-0 coefficient invariant under any deformation; higher orders not &
\lm{LambdaMomentHierarchy.order0\_deformation\_invariant}, \lm{only\_count\_touches\_lambda} & \Kernel\\
magnitude verdict: only count statistics move $\Lam$ &
\lm{LambdaMagnitudeCapstone.lambda\_only\_count\_can\_move\_order0} & \Kernel\\
finite unimodular trade / multiplier / gauge shift &
\lm{LambdaUnimodular.multiplier\_field\_equation}, \lm{vacuum\_shift\_is\_gauge} & \Kernel\\
everpresent scaling $\Lam_{\mathrm{rms}}=1/\sqrt V$ &
\lm{...EverpresentLambdaScaling.everpresentLambda\_rms\_eq\_inv\_sqrt\_volume} & \Kernel\\
routed through the native edge count &
\lm{LambdaEdgeCount.lambda\_rms\_eq\_inv\_sqrt\_count}, \lm{edgecount\_extensive} & \Kernel\\
$\Lam$ as susceptibility; Bernoulli bound $\mathrm{Var}(N)\le\langle N\rangle$ &
\lm{LambdaSusceptibility.var\_count}, \lm{bernoulli\_bound} & \Kernel\\
Poisson vs hyperuniform dichotomy on finite witnesses &
\lm{LambdaCountDichotomy.free\_extensive}, \lm{hard\_subextensive} & \Kernel\\
exponent law $\lm{lamExp}(\alpha)=\alpha/2-1$; $\lm{lamExp}(1)=-1/2$; fork iff &
\lm{LambdaExponentFork.lamExp\_closed}, \lm{everpresent\_value}, \lm{fork\_iff} & \Kernel\\
finite $\Lam$--volume conjugacy \& support uncertainty &
\lm{LambdaConjugacy.support\_uncertainty}, \lm{delta\_saturates} & \Kernel\\
$1/\sqrt V\sim10^{-122}$ for the present horizon & Ahmed--Dodelson--Greene--Sorkin & \Import\\
$\Lam$--$V$ conjugacy (unimodular gravity) & Henneaux--Teitelboim; Unruh & \Import\\
$\Lam$ $=$ order-0 spectral-action term & Chamseddine--Connes--Marcolli & \Import\\
which count is conjugate; kill on $\alpha<1$ & — & \Conj\\
value $10^{-122}$, sign, stochastic dynamics & — & outside\\
\bottomrule
\end{tabular}
\end{center}

\section*{References (verify before any manuscript citation) \Import}
\begin{enumerate}[leftmargin=1.6em,itemsep=0pt]
\item R.~Sorkin, ``Is the cosmological `constant' a nonlocal quantum residue of
  discreteness of the causal set type?'' (2007).
\item M.~Ahmed, S.~Dodelson, P.~Greene, R.~Sorkin, ``Everpresent Lambda,''
  Phys.\ Rev.\ D 69, 103523 (2004), arXiv:astro-ph/0209274.
\item T.~Jacobson, ``Thermodynamics of spacetime,'' arXiv:gr-qc/9504004;
  entanglement-equilibrium version arXiv:1505.04753.
\item A.~Chamseddine, A.~Connes, M.~Marcolli, ``Gravity and the standard model
  with neutrino mixing,'' arXiv:hep-th/0610241.
\item ``Bratteli networks and the spectral action on quivers,'' arXiv:2401.03705.
\item Unimodular gravity / $\Lam$--$V$ conjugacy: Henneaux--Teitelboim; Unruh.
\item S.~Torquato, F.~Stillinger, hyperuniformity; Martin--Yalcin,
  Stillinger--Lovett (Coulomb charge-fluctuation sum rules).
\item Planck 2018, arXiv:1807.06209; DESI DR2, arXiv:2503.14738;
  Zwane--Afshordi--Sorkin; ``Aspects-II'' everpresent-fluctuation bound.
\end{enumerate}

\end{document}
